\documentclass[a4paper,12pt]{article}
\usepackage[utf8]{inputenc}
\usepackage[T1]{fontenc}
\usepackage[ngerman]{babel}
\usepackage{amsmath,amssymb}
\usepackage{booktabs}
\usepackage{enumitem}
\usepackage[margin=2.5cm]{geometry}
\usepackage{listings}
\usepackage{xcolor}
\usepackage{tikz}
\lstset{basicstyle=\ttfamily\small, keywordstyle=\color{blue!70!black},
        commentstyle=\color{green!40!black}, stringstyle=\color{orange!80!black},
        numbers=left, numberstyle=\tiny, frame=single, breaklines=true,
        showstringspaces=false, language=Python}

\begin{document}

\begin{center}
{\LARGE\bfseries \"Ubungsblatt 4 -- L\"osungen}\\[0.5em]
{\large VC Logik}
\end{center}

\bigskip

%% ============================================================
%% AUFGABE 1
%% ============================================================
\subsection*{Aufgabe 1 \textnormal{(4 Punkte)} -- Sudoku-Formel}

\textbf{Angabe.} Entwerfen Sie eine aussagenlogische Formel, die ein
$4 \times 4$ Sudoku mit $2 \times 2$ Regionen repr\"asentiert (jedes
Modell = eine L\"osung). Jede Ziffer $1$--$4$ soll in jeder Zeile,
Spalte und Region h\"ochstens einmal vorkommen.

\medskip
\textbf{Spielregeln.}
\begin{itemize}[nosep]
\item Das Spielfeld ist ein $4 \times 4$-Gitter mit 16~Feldern.
\item Es ist in vier $2 \times 2$-Regionen unterteilt
      (links oben, rechts oben, links unten, rechts unten).
\item Jedes Feld muss mit genau einer Ziffer aus $\{1,2,3,4\}$
      bef\"ullt werden.
\item In jeder \textbf{Zeile} darf jede Ziffer h\"ochstens einmal
      vorkommen.
\item In jeder \textbf{Spalte} darf jede Ziffer h\"ochstens einmal
      vorkommen.
\item In jeder \textbf{Region} darf jede Ziffer h\"ochstens einmal
      vorkommen.
\end{itemize}
Da es 4~Ziffern und 4~Felder pro Zeile/Spalte/Region gibt, bedeutet
``h\"ochstens einmal'' hier sogar ``genau einmal''.

\medskip
\textbf{Beispiel eines Sudokus:}

\begin{center}
\begin{tikzpicture}[scale=0.9]
  \draw[thin] (0,0) grid (4,4);
  \draw[ultra thick] (0,0) rectangle (4,4);
  \draw[ultra thick] (2,0) -- (2,4);
  \draw[ultra thick] (0,2) -- (4,2);
  \node at (3.5,3.5) {\large 4};
  \node at (0.5,1.5) {\large 2};
  \node at (3.5,1.5) {\large 3};
  \node at (0.5,0.5) {\large 4};
  \node at (2.5,0.5) {\large 1};
  \node at (3.5,0.5) {\large 2};
\end{tikzpicture}
\end{center}
Die leeren Felder sollen so bef\"ullt werden, dass alle Regeln
erf\"ullt sind.

\medskip
\textbf{Schritt 1: Variablen einf\"uhren.}
F\"ur $i,j,n \in \{1,2,3,4\}$ definiere
\[
x_{i,j,n} \;:=\; \text{``Feld $(i,j)$ enth\"alt Ziffer $n$''}.
\]
Das sind $4 \cdot 4 \cdot 4 = 64$ Variablen.

\medskip
\textbf{Schritt 2: Regionen definieren.}
Die vier $2 \times 2$-Regionen:
\begin{align*}
R_1 &= \{(1,1),(1,2),(2,1),(2,2)\} \quad \text{(links oben)} \\
R_2 &= \{(1,3),(1,4),(2,3),(2,4)\} \quad \text{(rechts oben)} \\
R_3 &= \{(3,1),(3,2),(4,1),(4,2)\} \quad \text{(links unten)} \\
R_4 &= \{(3,3),(3,4),(4,3),(4,4)\} \quad \text{(rechts unten)}
\end{align*}

\medskip
\textbf{Schritt 3: Regeln formulieren.}
F\"unf Bedingungen als Teilformeln; die Gesamtformel ist ihre
Konjunktion.

\smallskip
\textit{Notation:} In den folgenden Formeln laufen alle Indizes von
$1$ bis $4$:
\begin{itemize}[nosep]
\item $i$: Zeilen-Index
\item $j$: Spalten-Index
\item $n, m$: Ziffern aus $\{1,2,3,4\}$
\item $r$: Regions-Index \;(w\"ahlt eine der 4 Regionen $R_1, \ldots, R_4$)
\end{itemize}
F\"ur die Vergleichs-Regeln (R3--R5) werden zwei verschiedene Felder
gebraucht. Die beiden Felder eines solchen Paars bekommen Indizes
mit Tiefstellung $_1$ und $_2$, also z.\,B.\ $i_1$ und $i_2$ f\"ur zwei
verschiedene Zeilen.

\smallskip
\textbf{(R1) Jedes Feld enth\"alt mindestens eine Ziffer:}
\[
R_1 \;\equiv\; \bigwedge_{i=1}^{4} \bigwedge_{j=1}^{4}
\bigl(x_{i,j,1} \vee x_{i,j,2} \vee x_{i,j,3} \vee x_{i,j,4}\bigr).
\]
In Worten: F\"ur jedes Feld $(i,j)$ muss mindestens eine der vier
Ziffern $1, 2, 3, 4$ dort gesetzt sein. $\Rightarrow$ 16~Klauseln
(eine pro Feld).

\smallskip
\textbf{(R2) Jedes Feld enth\"alt h\"ochstens eine Ziffer:}
\[
R_2 \;\equiv\; \bigwedge_{i=1}^{4} \bigwedge_{j=1}^{4}
\bigwedge_{1 \le n < m \le 4}
\neg\bigl(x_{i,j,n} \wedge x_{i,j,m}\bigr).
\]
In Worten: F\"ur jedes Feld $(i,j)$ und jedes Ziffern-Paar $n \ne m$
d\"urfen nicht beide gleichzeitig in diesem Feld stehen.
$\Rightarrow$ $16 \cdot 6 = 96$ Klauseln.

\smallskip
\textbf{(R3) Keine Ziffer zweimal in derselben Zeile:}
\[
R_3 \;\equiv\; \bigwedge_{i=1}^{4} \bigwedge_{n=1}^{4}
\bigwedge_{1 \le j_1 < j_2 \le 4}
\neg\bigl(x_{i,j_1,n} \wedge x_{i,j_2,n}\bigr).
\]
In Worten: F\"ur jede Zeile $i$, jede Ziffer $n$ und je zwei
verschiedene Spalten $j_1 < j_2$ d\"urfen die Felder $(i, j_1)$ und
$(i, j_2)$ nicht beide die Ziffer $n$ enthalten.
$\Rightarrow$ $4 \cdot 4 \cdot 6 = 96$ Klauseln.

\smallskip
\textbf{(R4) Keine Ziffer zweimal in derselben Spalte:}
\[
R_4 \;\equiv\; \bigwedge_{j=1}^{4} \bigwedge_{n=1}^{4}
\bigwedge_{1 \le i_1 < i_2 \le 4}
\neg\bigl(x_{i_1,j,n} \wedge x_{i_2,j,n}\bigr).
\]
In Worten: F\"ur jede Spalte $j$, jede Ziffer $n$ und je zwei
verschiedene Zeilen $i_1 < i_2$ d\"urfen die Felder $(i_1, j)$ und
$(i_2, j)$ nicht beide die Ziffer $n$ enthalten.
$\Rightarrow$ $4 \cdot 4 \cdot 6 = 96$ Klauseln.

\smallskip
\textbf{(R5) Keine Ziffer zweimal in derselben Region:}
\[
R_5 \;\equiv\; \bigwedge_{r=1}^{4} \bigwedge_{n=1}^{4}
\bigwedge_{\substack{(i_1,j_1),\,(i_2,j_2) \in R_r \\ (i_1,j_1) \neq (i_2,j_2)}}
\neg\bigl(x_{i_1,j_1,n} \wedge x_{i_2,j_2,n}\bigr).
\]
In Worten: F\"ur jede Region $R_r$, jede Ziffer $n$ und je zwei
verschiedene Felder $(i_1, j_1), (i_2, j_2)$ in dieser Region
d\"urfen die beiden Felder nicht beide die Ziffer $n$ enthalten.
$\Rightarrow$ $4 \cdot 4 \cdot 6 = 96$ Klauseln.

\bigskip
\noindent\fbox{\parbox{\dimexpr\textwidth-2\fboxsep-2\fboxrule}{%
\textbf{Ergebnis Aufgabe~1.}\\[2pt]
\[
\Phi \;\equiv\; R_1 \wedge R_2 \wedge R_3 \wedge R_4 \wedge R_5.
\]
Insgesamt $16 + 96 + 96 + 96 + 96 = 400$ Klauseln \"uber
$64$ Variablen. Jedes Modell von $\Phi$ entspricht einer g\"ultigen
Sudoku-L\"osung.
}}

\newpage

%% ============================================================
%% AUFGABE 2
%% ============================================================
\subsection*{Aufgabe 2 \textnormal{(2+2 Punkte)} -- Schaltkreis}

\textbf{Angabe.} Stelle den gegebenen Schaltkreis als aussagen"-logische
Formel dar, deren Modelle genau die zul\"assigen Konfigurationen der
Inputs/Outputs $A, B, C, D, E, F$ sind.
Notation: $[>k]$ gibt $1$ aus, wenn die Input-Summe $> k$ ist;
$[=k]$ gibt $1$ aus, wenn die Input-Summe $= k$ ist.

\medskip
\textbf{Schritt 1: Verkabelung ablesen.}
Aus dem Schaltkreis liest man folgende Struktur ab (mit
Hilfs-Outputs $o_1, o_2, o_3$ der linken Gates):

\begin{center}
\renewcommand{\arraystretch}{1.3}
\begin{tabular}{l|l|l|l}
\toprule
Gate & Inputs & Bedingung & Output \\
\midrule
$[>1]$ (oben links) & $A, B$ & $A+B > 1$ & $o_1$ \\
$[=1]$ (mitte links) & $A, C$ & $A+C = 1$ & $o_2$ \\
$[>0]$ (unten links) & $B, C$ & $B+C > 0$ & $o_3$ \\
$[=0]$ (oben rechts) & $o_1$ & $o_1 = 0$ & $D$ \\
$[>2]$ (mitte rechts) & $o_1, o_2, o_3$ & Summe $> 2$ & $E$ \\
$[>1]$ (unten rechts) & $o_2, o_3$ & $o_2 + o_3 > 1$ & $F$ \\
\bottomrule
\end{tabular}
\end{center}

\medskip
\textbf{Schritt 2: Gate-Bedingungen als Bool'sche Ausdr\"ucke.}
F\"ur bin\"are Inputs gilt:
\begin{align*}
A + B > 1 &\;\Leftrightarrow\; A \wedge B \\
A + C = 1 &\;\Leftrightarrow\; (A \wedge \neg C) \vee (\neg A \wedge C) \\
B + C > 0 &\;\Leftrightarrow\; B \vee C \\
o_1 = 0 &\;\Leftrightarrow\; \neg o_1 \\
o_1 + o_2 + o_3 > 2 &\;\Leftrightarrow\; o_1 \wedge o_2 \wedge o_3 \\
o_2 + o_3 > 1 &\;\Leftrightarrow\; o_2 \wedge o_3
\end{align*}

\medskip
\textbf{Schritt 3: Gesamtformel.}
F\"ur jeden Draht/Output wird eine \"Aquivalenz geschrieben und
alles mit $\wedge$ verkn\"upft:
\begin{align*}
\Phi \;\equiv\; & (o_1 \leftrightarrow (A \wedge B)) \wedge {} \\
& (o_2 \leftrightarrow ((A \wedge \neg C) \vee (\neg A \wedge C))) \wedge {} \\
& (o_3 \leftrightarrow (B \vee C)) \wedge {} \\
& (D \leftrightarrow \neg o_1) \wedge {} \\
& (E \leftrightarrow (o_1 \wedge o_2 \wedge o_3)) \wedge {} \\
& (F \leftrightarrow (o_2 \wedge o_3)).
\end{align*}
Die Modelle von $\Phi$ (projiziert auf $A,B,C,D,E,F$) sind genau
die zul\"assigen Konfigurationen.

\medskip
{\small
\textbf{Anmerkung.} Da $o_1, o_2, o_3$ Hilfsvariablen sind, lassen
sie sich auch eliminieren und direkt durch die folgende
\"aquivalente Formel ersetzen:
\begin{align*}
\Phi' \;\equiv\; & (D \leftrightarrow \neg(A \wedge B)) \wedge {} \\
& (E \leftrightarrow ((A \wedge B) \wedge ((A \wedge \neg C) \vee (\neg A \wedge C)) \wedge (B \vee C))) \wedge {} \\
& (F \leftrightarrow (((A \wedge \neg C) \vee (\neg A \wedge C)) \wedge (B \vee C))).
\end{align*}
Diese Form ist logisch \"aquivalent bez\"uglich der sichtbaren
Variablen $A,B,C,D,E,F$.
}

\medskip
\textbf{Schritt 4: Teilaufgaben (a), (b), (c).}
Alle drei Fragen lassen sich mit einem Erf\"ullbarkeitstester auf
$\Phi$ beantworten:

\begin{center}
\renewcommand{\arraystretch}{1.4}
\begin{tabular}{c|l|l}
\toprule
 & \textbf{Frage} & \textbf{Dem Tester geben} \\
\midrule
(a) & Kann $F = 1$ sein?
    & $\Phi \wedge F$
    \quad erf\"ullbar $\Rightarrow$ ja \\
(b) & Kann $E = A$ sein?
    & $\Phi \wedge (E \leftrightarrow A)$
    \quad erf\"ullbar $\Rightarrow$ ja \\
(c) & Gilt $F = 1 \Leftrightarrow D = 0$?
    & $\Phi \wedge \neg(F \leftrightarrow \neg D)$
    \quad unerf\"ullbar $\Rightarrow$ ja \\
\bottomrule
\end{tabular}
\end{center}

\medskip
\textbf{Konkrete Antworten (durch explizite Belegungen):}

\smallskip
\textbf{(a) Ja, $F = 1$ ist m\"oglich.}\\
W\"ahle $A = 1, B = 1, C = 0$. Dann:
$o_2 = (A \wedge \neg C) \vee (\neg A \wedge C) = (1 \wedge 1) \vee (0 \wedge 0) = 1$,
$o_3 = B \vee C = 1 \vee 0 = 1$,
$F = o_2 \wedge o_3 = 1 \wedge 1 = 1$. \checkmark

\smallskip
\textbf{(b) Ja, $E = A$ ist m\"oglich.}\\
Trivial f\"ur $A = 0$: w\"ahle $A = B = C = 0$, dann
$E = 0 \wedge \ldots = 0 = A$. \checkmark\\
Auch f\"ur $A = 1$: w\"ahle $A = 1, B = 1, C = 0$. Dann
$o_1 = 1$, $o_2 = 1$, $o_3 = 1$, also $E = 1 = A$. \checkmark

\smallskip
\textbf{(c) Nein, $F = 1 \Leftrightarrow D = 0$ gilt NICHT.}\\
Gegenbeispiel: $A = B = C = 1$. Dann:
$o_1 = A \wedge B = 1$, also $D = \neg o_1 = 0$.
$o_2 = (A \wedge \neg C) \vee (\neg A \wedge C) = 0$,
$o_3 = B \vee C = 1$, $F = o_2 \wedge o_3 = 0$.
Hier ist $D = 0$, aber $F = 0 \neq 1$. Also nicht \"aquivalent.

\newpage

%% ============================================================
%% AUFGABE 3
%% ============================================================
\subsection*{Aufgabe 3 \textnormal{(4 Punkte)} -- Bishop Independence
Problem}

\textbf{Angabe.} Schreiben Sie ein Skript, das Formeln produziert,
die das Bishop Independence Problem l\"osen (analog zu Rook/Queen
Independence Problem).

\medskip
\textbf{Schritt 1: Problem verstehen.}
Auf einem $n \times n$ Schachbrett: Platziere L\"aufer so, dass
keine zwei einander schlagen k\"onnen. L\"aufer bewegen sich
diagonal, d.\,h.\ zwei L\"aufer schlagen sich genau dann, wenn sie
auf derselben Diagonale stehen.

Diagonalen:
\begin{itemize}[nosep]
\item Hauptdiagonalen (von links oben nach rechts unten): $i - j$
      ist konstant, $i - j \in \{-(n-1), \ldots, n-1\}$ --
      insgesamt $2n - 1$ St\"uck.
\item Gegendiagonalen (von rechts oben nach links unten): $i + j$
      ist konstant, $i + j \in \{2, \ldots, 2n\}$ -- ebenfalls
      $2n - 1$ St\"uck.
\end{itemize}

\medskip
\textbf{Schritt 2: Variablen und Constraints.}
F\"ur jedes Feld $(i,j)$ mit $1 \le i,j \le n$ definiere
\[
b_{i,j} \;:=\; \text{``Auf Feld $(i,j)$ steht ein L\"aufer''}.
\]
F\"ur jede Diagonale $D$ (mit mindestens zwei Feldern) und jedes
Paar $(i_1,j_1), (i_2,j_2) \in D$ mit $(i_1,j_1) \neq (i_2,j_2)$:
\[
\neg\bigl(b_{i_1,j_1} \wedge b_{i_2,j_2}\bigr).
\]
Die Gesamtformel ist die Konjunktion all dieser Klauseln.

\medskip
\textbf{Schritt 3: Python-Skript.}

\begin{lstlisting}
def bishop_independence(n):
    """
    Erzeugt die aussagenlogische Formel (als Liste von Klauseln)
    fuer das Bishop Independence Problem auf einem n x n
    Schachbrett. Jede Klausel hat die Form
    (not b_{i1,j1} or not b_{i2,j2}) und verbietet zwei Laeufer
    auf derselben Diagonale.
    """
    # Alle Diagonalen sammeln
    diagonals = {}
    for i in range(1, n+1):
        for j in range(1, n+1):
            # Hauptdiagonale: i - j ist konstant
            diagonals.setdefault(('haupt', i-j), []).append((i, j))
            # Gegendiagonale: i + j ist konstant
            diagonals.setdefault(('gegen', i+j), []).append((i, j))

    # Klauseln erzeugen: fuer jede Diagonale mit >= 2 Feldern,
    # fuer jedes Paar ein at-most-one-Constraint
    clauses = []
    for feld_liste in diagonals.values():
        if len(feld_liste) < 2:
            continue
        for p in range(len(feld_liste)):
            for q in range(p+1, len(feld_liste)):
                (i1, j1) = feld_liste[p]
                (i2, j2) = feld_liste[q]
                klausel = f"(NOT b_{{{i1},{j1}}} OR NOT b_{{{i2},{j2}}})"
                clauses.append(klausel)
    return clauses


def formel_ausgabe(clauses):
    """Verbindet alle Klauseln mit AND zur Gesamtformel."""
    return " AND ".join(clauses)


if __name__ == "__main__":
    n = 4
    klauseln = bishop_independence(n)
    print(f"Anzahl Klauseln fuer n={n}: {len(klauseln)}")
    print(formel_ausgabe(klauseln))
\end{lstlisting}

\medskip
\textbf{Schritt 4: Beispiel $n = 4$.}
F\"ur ein $4 \times 4$-Brett:
\begin{itemize}[nosep]
\item Hauptdiagonalen: $7$ (f\"ur $i - j \in \{-3,-2,-1,0,1,2,3\}$).
      Gr\"o\ss en: $1, 2, 3, 4, 3, 2, 1$.
\item Gegendiagonalen: $7$ (f\"ur $i + j \in \{2,3,4,5,6,7,8\}$).
      Gr\"o\ss en: $1, 2, 3, 4, 3, 2, 1$.
\end{itemize}
Anzahl der Paare pro Diagonale: $\binom{k}{2}$ f\"ur Gr\"o\ss e~$k$.
Also $\binom{2}{2} + \binom{3}{2} + \binom{4}{2} + \binom{3}{2}
+ \binom{2}{2} = 1 + 3 + 6 + 3 + 1 = 14$ pro Richtung.
Insgesamt $14 + 14 = 28$ Klauseln.

\bigskip
\noindent\fbox{\parbox{\dimexpr\textwidth-2\fboxsep-2\fboxrule}{%
\textbf{Ergebnis Aufgabe~3.}\\[2pt]
Das Skript produziert f\"ur jedes $n$ eine Formel in CNF, die nur
die Independence-Constraints enth\"alt. Modelle = g\"ultige
L\"aufer-Konfigurationen (ohne gegenseitiges Schlagen). Triviale
Modelle (leeres Brett) sind eingeschlossen; wird eine Mindestanzahl
an L\"aufern gefordert, muss das als zus\"atzlicher Constraint
erg\"anzt werden.
}}

\end{document}
